\documentclass[aspectratio=169,11pt]{beamer}
% Shared CMPSC 480 Beamer preamble.
\usepackage[T1]{fontenc}
\usepackage[utf8]{inputenc}
\usepackage{lmodern}
\usepackage{amsmath,amssymb,mathtools}
\usepackage{booktabs,array,tabularx}
\usepackage{xcolor}
\usepackage{listings}
\usepackage{tikz}
\usetikzlibrary{arrows.meta,positioning,shapes.geometric}

% Ignore only negligible TeX box diagnostics that are below visual tolerance.
% Larger layout defects still appear in the build log and in rendered QA.
\vfuzz=3pt
\hbadness=2000

\definecolor{courseink}{HTML}{111111}
\definecolor{coursegray}{HTML}{EDEDED}
\definecolor{courserule}{HTML}{B8BCC4}
\definecolor{courseblue}{HTML}{3D8DFF}
\definecolor{courselightblue}{HTML}{D0EDFA}
\definecolor{coursemuted}{HTML}{5C6470}
\definecolor{coursegreen}{HTML}{0B7A53}
\definecolor{coursered}{HTML}{B3261E}

\usetheme{default}
\usefonttheme{professionalfonts}
\setbeamertemplate{navigation symbols}{}
\setbeamersize{text margin left=0.65cm,text margin right=0.65cm}

\setbeamercolor{normal text}{fg=courseink,bg=white}
\setbeamercolor{frametitle}{fg=courseink,bg=white}
\setbeamercolor{title}{fg=courseink,bg=white}
\setbeamercolor{subtitle}{fg=coursemuted,bg=white}
\setbeamercolor{structure}{fg=courseblue}
\setbeamercolor{alerted text}{fg=coursered}
\setbeamercolor{block title}{fg=courseink,bg=courselightblue}
\setbeamercolor{block body}{fg=courseink,bg=coursegray}
\setbeamercolor{example text}{fg=coursegreen}

\setbeamerfont{title}{size=\fontsize{25}{29}\selectfont,series=\bfseries}
\setbeamerfont{subtitle}{size=\fontsize{13}{16}\selectfont}
\setbeamerfont{frametitle}{size=\fontsize{19}{22}\selectfont,series=\bfseries}
\setbeamerfont{normal text}{size=\fontsize{11}{14}\selectfont}
\setbeamerfont{footline}{size=\fontsize{7.5}{9}\selectfont}

\setbeamertemplate{frametitle}{%
  \nointerlineskip
  % Use content-driven height so a long assessment title can wrap without
  % being clipped by a fixed-height title bar.
  \begin{beamercolorbox}[wd=\paperwidth,sep=0.17cm,leftskip=0.48cm,rightskip=0.48cm]{frametitle}
    \insertframetitle
  \end{beamercolorbox}
  \vspace{-0.08cm}
  {\color{courserule}\rule{\textwidth}{0.35pt}}
}

\setbeamertemplate{footline}{%
  \leavevmode
  \begin{beamercolorbox}[wd=\paperwidth,ht=2.6ex,dp=1.1ex,leftskip=.65cm,rightskip=.65cm]{author in head/foot}
    \color{coursemuted}\insertshorttitle\hfill\insertframenumber/\inserttotalframenumber
  \end{beamercolorbox}
}

\setbeamertemplate{itemize item}{\color{courseblue}\small$\blacktriangleright$}
\setbeamertemplate{itemize subitem}{\color{courseblue}\scriptsize$\bullet$}
\setbeamertemplate{enumerate item}{\color{courseblue}\insertenumlabel.}

\lstdefinestyle{coursepython}{
  language=Python,
  basicstyle=\ttfamily\fontsize{8.5}{10}\selectfont,
  keywordstyle=\color{courseblue}\bfseries,
  commentstyle=\color{coursemuted}\itshape,
  stringstyle=\color{coursegreen},
  showstringspaces=false,
  columns=fullflexible,
  keepspaces=true,
  frame=single,
  rulecolor=\color{courserule},
  backgroundcolor=\color{coursegray},
  breaklines=true,
  tabsize=4,
  xleftmargin=0.15cm,
  xrightmargin=0.15cm,
  framexleftmargin=0.15cm,
  framexrightmargin=0.15cm
}
\lstset{style=coursepython}

\newcommand{\points}[1]{\hfill{\color{courseblue}\textbf{[#1 points]}}}
\newcommand{\deliverable}[1]{\textbf{\color{courseblue}#1}}
\newcommand{\code}[1]{\texttt{#1}}
\newcommand{\keyidea}[1]{%
  \begin{beamercolorbox}[rounded=false,sep=0.55em,wd=\linewidth]{block body}
    \textbf{Key idea.} #1
  \end{beamercolorbox}
}
\newcommand{\answerlabel}{\textbf{\color{coursegreen}Answer.}\ }
\newcommand{\gradinglabel}{\textbf{\color{courseblue}Grading guidance.}\ }

\newenvironment{questionblock}[1]
  {\begin{block}{#1}}
  {\end{block}}

\newenvironment{solutionblock}[1]
  {\begin{block}{#1}}
  {\end{block}}

\AtBeginSection[]{
  \begin{frame}
    \vfill
    {\usebeamerfont{title}\color{courseink}\insertsectionhead\par}
    \vspace{0.4cm}
    {\color{courseblue}\rule{0.32\paperwidth}{3pt}}
    \vfill
  \end{frame}
}


% CMPSC480_TOTAL_POINTS: 20
% CMPSC480_ITEM_IDS: 1,2,3,4,5
\title[Module Project B Solutions]{Week 7 Module Project B: Q-Learning with Bayesian Belief Tracking}
\subtitle{Acting under partial observability with normalized posterior state}
\author{CMPSC 480 - Instructor Solution Deck}
\date{}

\begin{document}


\begin{frame}
  \titlepage
\vspace{-0.35cm}
\begin{center}
\color{coursered}\bfseries Instructor material - do not distribute with the question deck
\end{center}
\end{frame}

\begin{frame}{Solution overview}
\begin{columns}[T,onlytextwidth]
\begin{column}{0.62\textwidth}
\Large This instructor deck provides a complete matched solution and grading guidance.

\vspace{0.35cm}
\normalsize
Coverage: Weeks 4-7: MDPs, Q-learning, approximation, Bayes nets, and inference
\end{column}
\begin{column}{0.33\textwidth}
\begin{block}{Points}20 total\end{block}
\begin{block}{Expected time}10-14 hours\end{block}
\begin{block}{Submission}Repository, evaluation report, and belief traces\end{block}
\end{column}
\end{columns}
\end{frame}

\begin{frame}{Instructor verification checklist}
\small
\begin{itemize}
  \item Use Python 3.11 and NumPy; do not use a library that implements the assessed algorithm.
  \item Keep data, model, optimization, and evaluation responsibilities behind explicit interfaces.
  \item Validate shapes, ranges, finite values, empty inputs, and terminal or degenerate cases.
  \item Use deterministic seeds and record every configuration used for reported evidence.
  \item Include focused tests and a README with exact reproduction commands.
\end{itemize}
\end{frame}

\begin{frame}{Solution 1.a - Belief model (2 points)}
\small
\textbf{Question focus.} Write the vectorized posterior update.

\vspace{0.25cm}
\answerlabel For each hypothesis h, compute unnormalized[h]=likelihood(observation|h,agent\_state)*prior[h], then divide by the sum. Validate a nonnegative finite prior that sums to one.
\end{frame}

\begin{frame}{Solution 1.b - Belief model (2 points)}
\small
\textbf{Question focus.} Hand compute prior [0.5,0.5] with likelihoods [0.8,0.2] for one observation.

\vspace{0.25cm}
\answerlabel Unnormalized masses are [0.4,0.1], normalization is 0.5, and posterior is [0.8,0.2].
\end{frame}

\begin{frame}{Solution 1.c - Belief model (1 points)}
\small
\textbf{Question focus.} Define behavior when the normalization constant is zero or nonfinite.

\vspace{0.25cm}
\answerlabel Raise a domain-specific impossible-evidence error or apply a documented fallback such as restoring the prior; never silently divide by zero or invent a posterior.
\end{frame}

\begin{frame}{Solution 2.a - RL agent (2 points)}
\small
\textbf{Question focus.} Define the observable agent state and prove it excludes the true hidden hazard.

\vspace{0.25cm}
\answerlabel The state may include position, local observation, and belief vector or its summary. A test inspects inputs and confirms the true hazard identifier is never passed to policy or update code.
\end{frame}

\begin{frame}{Solution 2.b - RL agent (2 points)}
\small
\textbf{Question focus.} Specify epsilon-greedy selection and the terminal Q update.

\vspace{0.25cm}
\answerlabel Choose uniformly among legal actions with probability epsilon, otherwise stable argmax. For terminal steps target equals immediate reward; for others add gamma times max next Q.
\end{frame}

\begin{frame}{Solution 3.a - Integration contract (2 points)}
\small
\textbf{Question focus.} State the within-step order from observation to action to learning.

\vspace{0.25cm}
\answerlabel Condition the prior on the new observation, build the policy state from the updated belief, choose a legal action, execute it, then apply the reward transition update under the documented convention.
\end{frame}

\begin{frame}{Solution 3.b - Integration contract (1 points)}
\small
\textbf{Question focus.} Test that reset restores the declared prior and clears episode-local state.

\vspace{0.25cm}
\answerlabel Two episodes after reset begin with identical beliefs and agent state under the same seed, with no posterior or history leaked from the previous episode.
\end{frame}

\begin{frame}{Solution 3.c - Integration contract (1 points)}
\small
\textbf{Question focus.} List every component needed to resume training exactly.

\vspace{0.25cm}
\answerlabel Save Q parameters or table, belief configuration, online belief if mid-episode, epsilon, counters, environment state if supported, and all random-generator states.
\end{frame}

\begin{frame}{Solution 4.a - Noise study and ablation (2 points)}
\small
\textbf{Question focus.} Define safety and performance metrics and run multiple seeds.

\vspace{0.25cm}
\answerlabel Report mean return, success, hazard rate, and belief accuracy or entropy with variability for identical evaluation episodes and epsilon zero.
\end{frame}

\begin{frame}{Solution 4.b - Noise study and ablation (2 points)}
\small
\textbf{Question focus.} Ablate the posterior input and interpret the difference.

\vspace{0.25cm}
\answerlabel Replace the belief with a fixed prior or observation-only state while holding all else fixed. Attribute only the controlled performance change to belief integration.
\end{frame}

\begin{frame}{Solution 5.a - Defensive evidence (2 points)}
\small
\textbf{Question focus.} Provide four tests and name the defect each exposes.

\vspace{0.25cm}
\answerlabel Use repeated identical evidence, impossible evidence, uninformative sensor likelihoods, and a terminal transition; each test states expected belief or target and its associated bug.
\end{frame}

\begin{frame}{Solution 5.b - Defensive evidence (1 points)}
\small
\textbf{Question focus.} Explain why a high-noise posterior should remain closer to the prior.

\vspace{0.25cm}
\answerlabel When likelihoods are similar across hypotheses, multiplying by them changes relative prior masses little; normalized posterior information gain is therefore small.
\end{frame}

\begin{frame}{Edge-case reasoning and expected behavior}
\small
\begin{itemize}
  \item Reject the prior.
  \item Raise or use the documented impossible-evidence fallback.
  \item Posterior equals prior.
  \item Return without bootstrapping.
  \item Fail the leakage test and remove the hidden value.
\end{itemize}
\end{frame}

\begin{frame}{Grading rubric - part 1}
\scriptsize
\begin{tabularx}{\textwidth}{>{\bfseries}p{0.28\textwidth} p{0.08\textwidth} X}
\toprule
Category & Points & Full-credit evidence \\
\midrule
Belief model & 5 & Bayes numerator, normalization, repeated evidence, and validation are correct. \\
RL agent & 4 & No hidden-state leakage; exploration, targets, and termination are correct. \\
Integration contract & 4 & Update order, state encoding, reset, and episode isolation are correct. \\
\bottomrule
\end{tabularx}
\end{frame}

\begin{frame}{Grading rubric - part 2}
\scriptsize
\begin{tabularx}{\textwidth}{>{\bfseries}p{0.28\textwidth} p{0.08\textwidth} X}
\toprule
Category & Points & Full-credit evidence \\
\midrule
Noise study and ablation & 4 & Noise, seeds, budgets, baselines, metrics, and uncertainty are reported. \\
Defensive evidence & 3 & Tests cover impossible evidence, high noise, ties, and normalization drift. \\
\bottomrule
\end{tabularx}
\end{frame}

\begin{frame}{Reflection exemplar}
\small
\answerlabel A strong response verifies the posterior against an oracle, then uses visitation, reward design, policy state, or Q-learning evidence to locate the remaining control defect.

\vspace{0.25cm}
\gradinglabel Award credit for a specific claim supported by technical evidence from the submitted work.
\end{frame}

\end{document}
